\chapter{Power System Energy Management}

The robot is powered from a LiPo 2S battery with a nominal voltage of 7.4V and a maximum 
capacity of 6000\,mAh. This battery serves as the sole power source for all onboard systems. 
Power is distributed through two dedicated voltage rails: one supplying the servo motor 
subsystem via two switching regulators, and one supplying the computing and control subsystem 
via a single switching regulator.

% ----------------------------------------------------------------
\section{Power System Overview}
% INSERT DIAGRAM HERE

The power system is designed around a single LiPo 2S battery supplying all onboard subsystems 
through two dedicated power rails. The first rail services the servo motor subsystem, where two 
XL4016 DC/DC buck converters each regulate the battery voltage down to 5.5\,V, one per servo 
motor driver. Each driver manages six servo motors, for a total of twelve across the robot, and 
is rated for a maximum throughput of 8\,A. The second rail services the computing and control 
subsystem, where a single XL4015 buck converter regulates the battery voltage down to 5.2\,V, 
supplying both the STM32 microcontroller and the Raspberry Pi 5 from a single output rated at 
5\,A maximum. No additional protection circuitry is present in the power path.

% ----------------------------------------------------------------
\section{Battery Architecture, Power Distribution, and Runtime Estimation}

\subsection{Battery Architecture}

The battery is a 2S LiPo pack with a nominal voltage of 7.4\,V and a fully charged voltage of 
8.4\,V. The pack is composed of parallel cells, yielding a total capacity of 6000\,mAh 
(44.4\,Wh). The parallel cell configuration increases the available capacity and current 
delivery capability while maintaining the 2S series voltage level.

Power is distributed from the battery terminal to two independent rails as summarised in 
Table~\ref{tab:power_rails}.

\begin{table}[h]
    \centering
    \caption{Power Rail Summary}
    \label{tab:power_rails}
    \begin{tabular}{|l|c|c|c|l|}
        \hline
        \textbf{Rail} & \textbf{Regulator} & \textbf{Output Voltage} & \textbf{Max Current} & \textbf{Loads} \\
        \hline
        Servo Rail 1  & XL4016 & 5.5\,V & 8\,A & 6 $\times$ Servo Motors \\
        Servo Rail 2  & XL4016 & 5.5\,V & 8\,A & 6 $\times$ Servo Motors \\
        Compute Rail  & XL4015 & 5.2\,V & 5\,A & Raspberry Pi 5, STM32 MCU \\
        \hline
    \end{tabular}
\end{table}

\subsection{Full-Load Power Distribution}

The full-load current draw for each subsystem is as follows:

\begin{itemize}
    \item Each servo motor draws approximately 1.2\,A at full load. With twelve motors split 
          evenly across two drivers, each driver supplies $6 \times 1.2 = 7.2$\,A.
    \item The Raspberry Pi 5 draws up to 5\,A.
    \item The STM32 microcontroller draws approximately 500\,mA.
    \item The compute rail total draw is therefore $5 + 0.5 = 5.5$\,A.
\end{itemize}

The total system current demand at full load is summarised in Table~\ref{tab:current_draw}.

\begin{table}[h]
    \centering
    \caption{Full-Load Current Draw by Subsystem}
    \label{tab:current_draw}
    \begin{tabular}{|l|c|c|}
        \hline
        \textbf{Subsystem} & \textbf{Current Draw (A)} & \textbf{Rail} \\
        \hline
        6 $\times$ Servo Motors (Driver 1) & 7.2  & Servo Rail 1 \\
        6 $\times$ Servo Motors (Driver 2) & 7.2  & Servo Rail 2 \\
        Raspberry Pi 5                     & 5.0  & Compute Rail \\
        STM32 MCU                          & 0.5  & Compute Rail \\
        \hline
        \textbf{Total}                     & \textbf{19.9} & --- \\
        \hline
    \end{tabular}
\end{table}

\subsection{Runtime Estimation}

The theoretical runtime at full load is estimated using:

\begin{equation}
    t = \frac{C}{I_{\text{total}}}
    \label{eq:runtime}
\end{equation}

\noindent where $C$ is the battery capacity in ampere-hours and $I_{\text{total}}$ is the total 
current draw in amperes. Substituting the values:

\begin{equation}
    t = \frac{6.0\,\text{Ah}}{19.9\,\text{A}} \approx 0.30\,\text{h} \approx 18\,\text{minutes}
\end{equation}

This figure represents a theoretical lower bound under simultaneous full actuation of all servo 
motors and peak compute load. In practice, not all servo motors will be at peak draw 
simultaneously, and the Raspberry Pi 5 will seldom sustain its maximum current draw, so 
operational runtime is expected to exceed this estimate under normal working conditions.

% ----------------------------------------------------------------
\section{Power Circuit Design and Load Path}

\subsection{Load Path}

Battery power flows from the LiPo pack through the main power distribution node, from which 
the three switching regulators are fed in parallel. Each XL4016 regulator drives its dedicated 
servo motor driver, and the XL4015 regulator drives the compute rail shared by the Raspberry 
Pi 5 and the STM32 MCU.

\subsection{Wiring}

All power wiring throughout the system uses 18\,AWG wire. The current-carrying capacity of 
18\,AWG copper wire is typically rated at approximately 16\,A in free air for chassis wiring, 
which is sufficient for the individual rail currents in this design. Table~\ref{tab:wiring} 
summarises the wire assignments and the corresponding maximum currents on each segment.

\begin{table}[h]
    \centering
    \caption{Wiring Summary}
    \label{tab:wiring}
    \begin{tabular}{|l|c|c|}
        \hline
        \textbf{Segment}                  & \textbf{Wire Gauge} & \textbf{Max Current (A)} \\
        \hline
        Battery to distribution node      & 18\,AWG & 19.9 \\
        Distribution node to XL4016 \#1   & 18\,AWG & 7.2  \\
        Distribution node to XL4016 \#2   & 18\,AWG & 7.2  \\
        Distribution node to XL4015       & 18\,AWG & 5.5  \\
        XL4016 \#1 to Servo Driver 1      & 18\,AWG & 7.2  \\
        XL4016 \#2 to Servo Driver 2      & 18\,AWG & 7.2  \\
        XL4015 to Raspberry Pi 5          & 18\,AWG & 5.0  \\
        XL4015 to STM32 MCU               & 18\,AWG & 0.5  \\
        \hline
    \end{tabular}
\end{table}

\noindent It should be noted that the battery-to-distribution segment carries the full system 
current of up to 19.9\,A at peak load. While 18\,AWG wire is within rated limits for this 
current, careful attention to wire length and thermal conditions is advised for this segment 
in particular to minimise resistive losses and heat generation.
